\section{Introduction}
Quantum programming tools expose circuits, operators, and execution interfaces,
whereas the maintainer of a classical application starts with program behavior.
The maintainer needs to determine \emph{where} a candidate computation occurs,
\emph{whether} a supported reformulation is justified, and \emph{how} the
surrounding software contract can be retained. These questions arise before
committing to a quantum implementation. Retaining the classical computation is
a legitimate outcome, including when a conditional quantum mapping can be described.

This problem resembles selective accelerator offloading: a computational kernel
is only one component of a replacement. Dependencies, marshaling, selection rules,
fallback, and externally visible behavior can dominate both correctness and cost.
A graph optimization routine, for example, may return the first optimum according
to a specified ordering. An optimizer that returns an arbitrary optimum can solve
the mathematical problem while violating the application contract. Similarly,
finding one satisfying assignment is not equivalent to returning a complete ordered
report, and a failed randomized search does not certify the absence of a solution.

Existing work provides important foundations. C2$|$Q$\rangle$ translates classical
specifications into quantum programs through encoding, deployment, and decoding
components~\cite{c2q}. Qiskit HumanEval and QuanBench evaluate quantum code
generation using explicit programming tasks and executable
checks~\cite{qhe,quanbench}. QuanBench+ also examines repair after execution
feedback~\cite{quanplus}. These directions motivate a complementary question:
can a fixed LLM reason more reliably about the obligations introduced when a
quantum mapping is proposed \emph{inside an existing classical program}?
We do not claim that feedback-based generation itself is new.

The central difficulty is connecting a recognized mathematical pattern to the
source program's obligations. Input validation, selection order, and intermediate
states can invalidate an otherwise plausible mapping. A local formula can pass
finite tests while certification and integration remain unresolved.

\bench{} supports studying these distinctions. Its development material adapts
traceable source tasks into classical programs with functionally relevant context.
The supported positive families are deliberately limited to unstructured search
and combinatorial optimization. Core-only, hinted-context, and unhinted-context
views permit controlled diagnosis. Private reference proposals and semantic
checks are kept separate from model-visible inputs. The current reference layer
is provisional; it is not described as independently validated ground truth.

Building on this foundation, we propose a workflow that attaches program facts
and contract obligations to model claims, checks supported representations, and
returns scoped feedback to the same model. Its central hypothesis is that
\emph{evidence tied to the source program can improve mapping design beyond the
benefit of another model call}. The relevant comparison therefore holds the
model fixed and varies the assistance. Cross-model runs assess the reach of a
finding; they do not replace this within-model comparison.

The intended paper makes three contributions: (1) a task and benchmark design
that separates candidate discovery, structural mapping, practical adoption, and
behavioral obligations; (2) a proposed integration of program facts, structured
claims, and semantic counterexamples into a fixed-model workflow; and (3) a
controlled empirical study of correctness, completion, failure modes, and cost.
Only the benchmark infrastructure and a bounded semantic-feedback prototype
currently exist. The integrated analysis component and the confirmatory study
remain proposed. This draft specifies the study and leaves all empirical result
cells blank, rather than supplying hypothetical improvements.

\figplaceholder{overview}{2.4cm}{Classical program and public contract $\rightarrow$
candidate region $\rightarrow$ structural mapping and obligation ledger
$\rightarrow$ conditional plan. Show program-fact and semantic-feedback loops
to one LLM; keep adoption, migration validation, and resource evidence separate.}
{Proposed workflow and evidence boundaries. The quantum component is a potential
target of a conditional plan; the workflow does not require quantum execution.}
